% For examples, Wang et al. showed that moderate scour depths markedly amplify seismic-demand fragilities \cite{Wang2014}; Guo and Chen introduced a life-cycle multi-hazard framework coupling stochastic scour progression with earthquake loading \cite{Guo2015}; Fan et al. demonstrated how scour, corrosion, and barge impact interact to shift failure from superstructure to foundations \cite{Fan2021}; Anisha et al. quantified that a $1.5-m$ scour depth can increase flood-induced failure probability by approximately $80\%$ \cite{Anisha2022}; and Karriqi et al. confirmed that scour-related stiffness loss modestly elevates seismic fragility when earthquake and flood hazards act sequentially \cite{Karriqi2024}.

% Among these efforts, although variations exist, they feature a similar approach to estimating structural collapse probabilistically. For a generic structure, the approach starts with defining a displacement-based capacity attribute, which is often taken as the yielding value of the structure, denoted as $\Delta_y$. Such capacity measurement can be obtained from a nonlinear model pushover analysis. Coressponding from the dislacement variable at the select location and direction, a displacement-based response-demand variable, denoted as $\Delta$, can be measured from simulating the brige structure subject to the hazardous loading, defined by an intensity measure $H$. Given such measurements from simulating the structure, a collapse event can be defined as a limit state. Often a ratio is defined, $u = {\Delta}{\Delta_y}$; then the limit state is expressed as $u > u_t$ with $u_t$ being a select threshold parameter. Through stochastic simulation of the structure considering loading and material uncertainties, researchers can then estimate the conditinoal probability of the event $P(u > u_t | H )$. For structures without varying scour profiles or boundary conditions at the foundtion level, the capacity variable $\Delta_y$ is a constant as an attribute of the structure. However, for scour-critical structures, $\Delta_y$ varies with the boundary condition at the foundation level due to the scour-profile variation.  

% To illustrate this, Figure~\ref{fig:figure1a} depicts three nonlinear pushover profiles considering a bridge model with variable scour profiles: the solid-curve with the yielding capacity pair of $(\Delta_{y0}, Q_{y0})$ indicates a non-scour situation, the dashed curve with $(\Delta_{y1}, Q_{y1})$ for a moderate scour case, and the dot-dashed curve with $(\Delta_{y2}, Q_{y2})$ for a severe scour case. Given the same predictve displacement response of $\Delta$ from the bridge model with three parameteric scour profiles, if the collapse potential index is defined using the non-scour capacity, $\frac{\Delta}{\Delta_{y0}}$, it implies the same collapse potential. In other words, the limit state defined as $u = \frac{\Delta}{\Delta_y} > u_t$, although commonly used in the existing efforts, is problematic. Furthermore, given the thress nonlinear pushover profiles in Figure~\ref{fig:figure1a}, it is straightforward that the collapse potential for scour-critical bridges needs to consider both displacement and strength capacity, namely, including the reduced strength $Q_y$ and $\Delta_y$; consequetially, the collapse limit state is defined. 

% \begin{figure}[htbp]
%     \centering
%     \begin{subfigure}[b]{0.48\textwidth}
%         \centering
%         \includegraphics[width=\linewidth]{Figures/Figure01a.pdf}
%         \caption{}
%         \label{fig:figure1a}
%     \end{subfigure}
%     \hfill
%     \begin{subfigure}[b]{0.48\textwidth}
%         \centering
%         \includegraphics[width=\linewidth]{Figures/Figure01b.pdf}
%         \caption{}
%         \label{fig:figure1b}
%     \end{subfigure}
%     \caption{(a) Illustration of inconsistency in the definition of collapse indices for scour-critical bridge structures; (b) proposed surrogate modeling and uncertainty quantification framework and methodological components.}
%     \label{fig:figure1}
% \end{figure}



% The diversity of limit-state metrics complicates cross-study comparisons and motivates the capacity-tuple adopted in the present work. Table~\ref{tab:lit_summary} (see appendix) summarizes key studies addressing scour and multi-hazard effects on bridge performance, highlighting methods, capacity metrics, uncertainty treatment, and remaining research gaps.

% \begin{table*}[htbp]
% \caption{Representative Studies on Scour and Multi-Hazard Effects on Bridge Performance}
% \label{tab:lit_summary}
% \centering
% \resizebox{\textwidth}{!}{%
% \begin{tabular}{|p{2.8cm}|p{2.4cm}|p{3.2cm}|p{2.8cm}|p{3.0cm}|}
% \hline
% \textbf{Study} & \textbf{Hazards Considered} & \textbf{Method} & \textbf{Capacity Metric(s)} & \textbf{Uncertainty/Remarks} \\
% \hline
% Prasad and Banerjee (2013) & Scour & Probabilistic fragility, nonlinear pushover & $\Delta_y$ & Scour as RV, basic fragility shift \\
% \hline
% Guo, Badroddin, and Chen (2019) & Scour & Empirical fragility fitting (inspection data) & $\Delta_y$ & Scour-dependent empirical model \\
% \hline
% Annad et al. (2023) & Scour & 3D FE w/ spring removal, case study & $\Delta_y$, $V_y$ & Soil-layered spring deactivation \\
% \hline
% Wang et al. (2014) & Scour + Seismic & Nonlinear pushover, probabilistic fragility & $\Delta_y$ & Scour raises column failure fragility \\
% \hline
% Guo and Chen (2015); Guo et al. (2016) & Scour + Seismic & Time-dependent lifecycle fragility & $\Delta_y$ & Lifecycle migration of fragility curves \\
% \hline
% Argyroudis and Mitoulis (2021) & Scour + Seismic & Copula fragility, parametric methods & $\Delta_y$ & Confirmed fragility elevation across typologies \\
% \hline
% Kameshwar and Padgett (2018) & Scour + Barge Impact & Probabilistic analysis of lateral loss & $\theta_y$, pile response & Pier scour > 1.5 diameters triggers collapse \\
% \hline
% Fan et al. (2021) & Scour + Corrosion + Barge & FE modeling, system interaction & $\mu_\phi$, pile hinge & Shift from deck to pile failure \\
% \hline
% Anisha et al. (2022) & Scour + Debris/Flood & Joint hazard modeling, fragility & $P_f$ under flood & 1.5m scour $\Rightarrow$ +80\% flood fragility \\
% \hline
% Karriqi et al. (2024) & Seismic $\rightarrow$ Scour (Seq) & Time-sequence pushover fragility & System-level fragility & Sequential loading increases fragility \\
% \hline
% \end{tabular}%
% }
% \end{table*}